% Part: turing-machines
% Chapter: machines-computations
% Section: representing-tms

\documentclass[../../../include/open-logic-section]{subfiles}

\begin{document}

\olfileid[hi]{tur}{mac}{rep}
\olsection{ट्यूरिंग मशीनों का निरूपण}

\begin{explain}
ट्यूरिंग मशीनों को \emph{अवस्था आरेखों} द्वारा दृश्य रूप में निरूपित किया
जा सकता है। ये आरेख तीरों से जुड़े अवस्था-कोष्ठों से बने होते हैं। जैसा
कि अपेक्षित है, प्रत्येक अवस्था-कोष्ठ मशीन की एक अवस्था को निरूपित करता
है। प्रत्येक तीर उस निर्देश को निरूपित करता है जिसे उस अवस्था से
निष्पादित किया जा सकता है; निर्देश का विशिष्ट विवरण संबंधित तीर के ऊपर
या नीचे लिखा होता है। निम्न मशीन पर विचार करें, जिसमें केवल दो आंतरिक
अवस्थाएँ $q_0$ और $q_1$, तथा एक निर्देश है:
\[
\begin{tikzpicture}[->,>=stealth',shorten >=1pt,auto,node distance=2.8cm,
                    semithick]
  \tikzstyle{every state}=[fill=none,draw=black,text=black]

  \node[initial,state]   (A)              {$q_0$};
  \node[state]   (B) [right of=A] {$q_1$};

  \path (A) edge  node {\TMtrans{\TMblank}{\TMstroke}{\TMright}} (B);
\end{tikzpicture}
\]
याद करें कि ट्यूरिंग मशीन में एक पठन-लेखन शीर्ष और उस पर लिखे आगत वाली
एक पट्टी होती है। इस निर्देश को ऐसे पढ़ सकते हैं: \emph{यदि
एक~$\TMblank$ को अवस्था $q_0$ में पढ़ रहे हों, तो~$\TMstroke$ लिखें, दाईं ओर जाएँ और
अवस्था $q_1$ में चले जाएँ}। यह उस संक्रमण फलन के समतुल्य है जो
$\tuple{q_0, \TMblank}$ को $\tuple{q_1, \TMstroke,
\TMright}$ पर भेजता
है।
\end{explain}

\begin{ex}
\emph{सम मशीन}: निम्न ट्यूरिंग मशीन तभी और केवल तभी रुकती है जब पट्टी पर
$\TMstroke$ प्रतीकों की संख्या सम हो (यह मानते हुए कि सभी $\TMstroke$
प्रतीक पट्टी के पहले $\TMblank$ से पहले आते हैं)।
\[
\begin{tikzpicture}[->,>=stealth',shorten >=1pt,auto,node distance=2.8cm,
                    semithick]
  \tikzstyle{every state}=[fill=none,draw=black,text=black]

  \node[initial,state]         (A)              {$q_0$};
  \node[state]         (B) [right of=A] {$q_1$};

  \path (A) edge [bend left] node {\TMtrans{\TMstroke}{\TMstroke}{\TMright}} (B)
        (B) edge [loop above] node {\TMtrans{\TMblank}{\TMblank}{\TMright}} (B)
            edge [bend left] node {\TMtrans{\TMstroke}{\TMstroke}{\TMright}} (A);
\end{tikzpicture}
\]

यह अवस्था आरेख निम्न संक्रमण फलन के अनुरूप है:
\begin{align*}
\delta(q_0, \TMstroke) & = \tuple{q_1, \TMstroke, \TMright},\\
\delta(q_1, \TMstroke) & = \tuple{q_0, \TMstroke, \TMright},\\
\delta(q_1, \TMblank)  & = \tuple{q_1, \TMblank, \TMright}
\end{align*}
\end{ex}

\begin{explain}
ऊपर की मशीन केवल तभी रुकती है जब आगत में आघात-प्रतीकों की संख्या सम हो।
अन्यथा मशीन (सैद्धांतिक रूप से) अनिश्चित काल तक चलती रहती है। किसी भी
मशीन और आगत के लिए निर्गत निर्धारित करने हेतु मशीन के
\emph{विन्यासों} का क्रम देखा जा सकता है। हम विन्यास की आकारिक परिभाषा
आगे देंगे। अभी हम सहज रूप से विन्यासों को ऐसे आरेखों की शृंखला मान सकते
हैं जो क्रिया के दौरान किसी भी समय मशीन की अवस्था दिखाते हैं। विन्यास
पट्टी की सामग्री, मशीन की अवस्था और पठन-लेखन शीर्ष का स्थान दिखाते हैं।

यदि सम मशीन को चार $\TMstroke$ के आगत से आरंभ किया जाए, तो उसके विन्यासों
का क्रम देखें। इस स्थिति में हम मशीन के रुकने की अपेक्षा करते हैं। इसके
बाद मशीन को तीन $\TMstroke$ के आगत पर चलाएँगे, जहाँ वह सदा चलती रहेगी।

मशीन अवस्था~$q_0$ में आरंभ होकर सबसे बाएँ~$\TMstroke$ को पढ़ती है। मशीन
के आरंभिक विन्यास को इस प्रकार निरूपित कर सकते हैं:
\[
\TMendtape \TMstroke_0 \TMstroke \TMstroke \TMstroke \TMblank \ldots
\]
ऊपर का विन्यास सीधा है। जैसा दिखाई देता है, मशीन आरंभिक अवस्था में
सबसे बाएँ~$\TMstroke$ को पढ़ती है। इसे पहले~$\TMstroke$ पर अवस्था
के नाम को अधोलिखित करके निरूपित किया गया है। इस बिंदु पर लागू निर्देश
$\delta(q_0, \TMstroke) = \tuple{q_1, \TMstroke, \TMright}$ है; अतः
मशीन पट्टी पर दाईं ओर जाती है और अवस्था~$q_1$ में बदल जाती है।
\[
\TMendtape \TMstroke \TMstroke_1 \TMstroke \TMstroke \TMblank \ldots
\]
चूँकि अब मशीन अवस्था~$q_1$ में एक~$\TMstroke$ पढ़ रही है, हमें निर्देश
$\delta(q_1, \TMstroke) = \tuple{q_0,
  \TMstroke, \TMright}$ का
``अनुसरण'' करना होगा। इससे निम्न विन्यास मिलता है:
\[
\TMendtape \TMstroke \TMstroke \TMstroke_0 \TMstroke \TMblank \ldots
\]
मशीन के आगे चलने पर नियम पुनः उसी क्रम में लागू होते हैं, जिससे अगले दो
विन्यास मिलते हैं:
\[
\TMendtape \TMstroke \TMstroke \TMstroke \TMstroke_1 \TMblank \ldots
\]
\[
\TMendtape \TMstroke \TMstroke \TMstroke \TMstroke \TMblank_0 \ldots
\]
अब मशीन अवस्था~$q_0$ में एक~$\TMblank$ पढ़ रही है। संक्रमण आरेख से सरलता
से दिखता है कि निष्पादित करने के लिए कोई निर्देश नहीं है; अतः मशीन रुक
गई है। इसका अर्थ है कि आगत स्वीकार कर लिया गया है।

अब मान लें कि मशीन को तीन $\TMstroke$ के आगत से आरंभ करते हैं। आरंभिक
कुछ विन्यास समान हैं, क्योंकि वही निर्देश निष्पादित होते हैं; अंतर केवल
पट्टी के आगत में थोड़ा-सा है:
\[
\TMendtape \TMstroke_0 \TMstroke \TMstroke \TMblank \ldots
\]
\[
\TMendtape \TMstroke \TMstroke_1 \TMstroke \TMblank \ldots
\]
\[
\TMendtape \TMstroke \TMstroke \TMstroke_0 \TMblank \ldots
\]
\[
\TMendtape \TMstroke \TMstroke \TMstroke \TMblank_1 \ldots
\]
अब मशीन सभी $\TMstroke$ से आगे निकल चुकी है और एक~$\TMblank$ को
अवस्था~$q_1$ में पढ़ रही है। आरेख के अनुसार यहाँ
$\delta(q_1, \TMblank) =\tuple{q_1,
\TMblank, \TMright}$ रूप का निर्देश
है। चूँकि पट्टी दाईं ओर अनिश्चित काल तक $\TMblank$ से भरी है, मशीन
अवस्था~$q_1$ में रहते और निरंतर दाईं ओर जाते हुए इस निर्देश को
\emph{सदा} निष्पादित करती रहेगी। मशीन कभी नहीं रुकेगी और आगत स्वीकार
नहीं करेगी।
\end{explain}

\begin{explain}
यह ध्यान रखना महत्त्वपूर्ण है कि सभी मशीनें नहीं रुकेंगी। यदि रुकने का
अर्थ यह है कि मशीन के पास निष्पादित करने के लिए कोई निर्देश शेष नहीं
रहा, तो प्रत्येक अवस्था पर प्रत्येक प्रतीक के लिए एक बाहर जाने वाला तीर
सुनिश्चित करके ऐसी मशीन बनाई जा सकती है जो कभी न रुके। एक
$\TMblank$ को अवस्था~$q_0$ में पढ़ने का निर्देश जोड़कर सम मशीन को अनिश्चित काल तक चलने वाली
मशीन में बदला जा सकता है।
\end{explain}

\begin{ex}
\[
\begin{tikzpicture}[->,>=stealth',shorten >=1pt,auto,node distance=2.8cm,
                    semithick]
  \tikzstyle{every state}=[fill=none,draw=black,text=black]

  \node[initial,state] (A)              {$q_0$};
  \node[state]         (B) [right of=A] {$q_1$};

  \path
  (A) edge [bend left] node {\TMtrans{\TMstroke}{\TMstroke}{\TMright}} (B)
      edge [loop above] node {\TMtrans{\TMblank}{\TMblank}{\TMright}} (A)
  (B) edge [loop above] node {\TMtrans{\TMblank}{\TMblank}{\TMright}} (B)
      edge [bend left] node {\TMtrans{\TMstroke}{\TMstroke}{\TMright}} (A);
\end{tikzpicture}
\]
\end{ex}

\begin{explain}
मशीन सारणियाँ ट्यूरिंग मशीनों के निरूपण की दूसरी रीति हैं। मशीन सारणी
में पट्टी की वर्णमाला $x$-अक्ष पर और मशीन की अवस्थाओं का समुच्चय
$y$-अक्ष पर दिखाया जाता है। सारणी के भीतर प्रत्येक अवस्था और प्रतीक के
प्रतिच्छेद पर निर्देश का शेष भाग---नई अवस्था, नया प्रतीक और गति की
दिशा---लिखा होता है। मशीन सारणी से यह निर्धारित करना सरल हो जाता है कि
मशीन किस अवस्था और किस प्रतीक पर रुकती है। सारणी का प्रत्येक रिक्त स्थान
मशीन के रुकने का संभावित बिंदु है। अवस्था आरेखों और निर्देश-समुच्चयों में
मशीन के रुकने के बिंदु सदैव तुरंत स्पष्ट नहीं होते; इसके विपरीत, मशीन
सारणी के रिक्त स्थान खोजकर सभी विराम-बिंदुओं को शीघ्र पहचाना जा सकता है।
\end{explain}

\begin{ex}
सम मशीन की मशीन सारणी यह है:
\[
\centering
\begin{tabular}{lllll}
\cline{2-4}
\multicolumn{1}{l|}{}      & \multicolumn{1}{c|}{$\TMblank$}                
& \multicolumn{1}{c|}{$\TMstroke$}     & \multicolumn{1}{c|}{$\TMendtape$}   &  \\ \cline{1-4}
\multicolumn{1}{|l|}{$q_0$} & \multicolumn{1}{l|}{}                          
& \multicolumn{1}{l|}{\TMtrans{\TMstroke}{q_1}{\TMright}} & 
\multicolumn{1}{l|}{\phantom{\TMtrans{\TMstroke}{q_1}{\TMright}}} &  \\ \cline{1-4}
\multicolumn{1}{|l|}{$q_1$} & \multicolumn{1}{l|}{\TMtrans{\TMblank}{q_1}{\TMright}} 
& \multicolumn{1}{l|}{\TMtrans{\TMstroke}{q_0}{\TMright}} & 
\multicolumn{1}{l|}{\phantom{\TMtrans{\TMstroke}{q_1}{\TMright}}} &  \\ \cline{1-4}
\end{tabular}
\]
जैसा दिखाई देता है, मशीन एक~$\TMblank$ को अवस्था~$q_0$ में पढ़ते समय रुकती
है।
\end{ex}

\begin{explain}
अब तक हमने केवल उन मशीनों पर विचार किया है जो आगत को पढ़ती और स्वीकार
करती हैं। किंतु ट्यूरिंग मशीनों में पढ़ने और लिखने दोनों की क्षमता होती
है। ऐसी मशीन का एक उदाहरण (यद्यपि उदाहरण बहुत अधिक हैं) \emph{द्विगुणक}
है। पट्टी पर~$n$ $\TMstroke$ के खंड से आरंभ करने पर द्विगुणक~$2n$
$\TMstroke$ का खंड निर्गत करता है।
\end{explain}

\begin{ex}
  \ollabel{ex:doubler} द्विगुणक मशीन बनाने से पहले समस्या हल करने की
  एक \emph{रणनीति} बनाना महत्त्वपूर्ण है। चूँकि मशीन (जैसा हमने उसे
  सूत्रबद्ध किया है) यह याद नहीं रख सकती कि उसने कितने $\TMstroke$ पढ़े
  हैं, इसलिए पट्टी के सभी $\TMstroke$ का लेखा रखने की रीति चाहिए। एक रीति
  यह है कि निर्गत को आगत से एक~$\TMblank$ द्वारा अलग कर दें। तब मशीन आगत
  से पहला $\TMstroke$ मिटा सकती है, शेष आगत को पार कर सकती है, एक
  $\TMblank$ छोड़ सकती है और दो नए $\TMstroke$ लिख सकती है। फिर मशीन
  लौटकर आगत में दूसरा $\TMstroke$ खोजेगी और उसे भी दुगुना करेगी। आगत के
  प्रत्येक एक $\TMstroke$ के लिए वह निर्गत के दो $\TMstroke$ लिखेगी।
  चलते-चलते आगत मिटाने से यह सुनिश्चित होता है कि कोई $\TMstroke$ छूटे
  नहीं और किसी को दो बार दुगुना न किया जाए। पूरा आगत मिट जाने पर पट्टी पर
  $2n$ $\TMstroke$ शेष रहेंगे। परिणामी ट्यूरिंग मशीन का अवस्था आरेख
  \olref{fig:doubler} में दिखाया गया है।
  \begin{figure}
\[
\begin{tikzpicture}[->,>=stealth',shorten >=1pt,auto,node distance=2.8cm,
                    semithick]
  \tikzstyle{every state}=[fill=none,draw=black,text=black]

  \node[initial,state] (1)              {$q_0$};
  \node[state]         (2) [right of=1] {$q_1$};
  \node[state]         (3) [right of=2] {$q_2$};
  \node[state]         (4) [below of=3] {$q_3$};
  \node[state]         (5) [left of=4]  {$q_4$};
  \node[state]         (6) [left of=5]  {$q_5$};

  \path (1) edge node {\TMtrans{\TMstroke}{\TMblank}{\TMright}} (2)
    (2) edge [loop above] node {\TMtrans{\TMstroke}{\TMstroke}{\TMright}} (2)
      edge node {\TMtrans{\TMblank}{\TMblank}{\TMright}} (3)
    (3) edge [loop above] node {\TMtrans{\TMstroke}{\TMstroke}{\TMright}} (3)
        edge node {\TMtrans{\TMblank}{\TMstroke}{\TMright}} (4)
    (4) edge [loop below] node {\TMtrans{\TMblank}{\TMstroke}{\TMleft}} (4)
        edge node {\TMtrans{\TMstroke}{\TMstroke}{\TMleft}} (5)
    (5) edge [loop below]  node {\TMtrans{\TMstroke}{\TMstroke}{\TMleft}} (5)
        edge              node {\TMtrans{\TMblank}{\TMblank}{\TMleft}} (6)
    (6) edge [loop below] node {\TMtrans{\TMstroke}{\TMstroke}{\TMleft}} (6)
        edge              node {\TMtrans{\TMblank}{\TMblank}{\TMright}} (1);
\end{tikzpicture}
\]
\caption{एक द्विगुणक मशीन}
\ollabel{fig:doubler}
\end{figure}
\end{ex}

\begin{prob}
कोई मनमाना आगत चुनें और \olref[tur][mac][rep]{ex:doubler} की द्विगुणक
मशीन के विन्यासों का क्रम निकालें।
\end{prob}

\begin{prob}
$\{\TMendtape,\TMblank, A, B\}$ वर्णमाला वाली ऐसी ट्यूरिंग मशीन की
अभिकल्पना करें जो $A$ और $B$ की प्रत्येक ऐसी अक्षर-शृंखला को स्वीकार
करे, अर्थात् उस पर रुके, जिसमें $A$ की संख्या $B$ की संख्या के बराबर हो
\emph{और} सभी $A$, सभी $B$ से पहले आएँ; तथा प्रत्येक ऐसी अक्षर-शृंखला
को अस्वीकार करे, अर्थात् उस पर न रुके, जिसमें $A$ की संख्या~$B$ की
संख्या के बराबर न हो अथवा सभी $A$, सभी~$B$ से पहले न आएँ। (उदाहरणतः,
मशीन को $AABB$ और $AAABBB$ स्वीकार करना चाहिए, किंतु $AAB$ और
$AABBAABB$ दोनों को अस्वीकार करना चाहिए।)
\end{prob}

\begin{prob}
$\{\TMendtape,\TMblank, A, B\}$ वर्णमाला वाली ऐसी ट्यूरिंग मशीन की
अभिकल्पना करें जो किसी भी अक्षर-शृंखला $\alpha$ को, जो $A$ और $B$ की हो, आगत के
रूप में लेकर उसकी प्रतिलिपि बनाए और $\alpha\alpha$ रूप का निर्गत दे।
(उदाहरणतः, आगत $ABBA$ का निर्गत $ABBAABBA$ होना चाहिए।)
\end{prob}

\begin{prob}
\emph{वर्णानुक्रम में?}: $\{\TMendtape,\TMblank, A, B\}$ वर्णमाला वाली
ऐसी ट्यूरिंग मशीन की अभिकल्पना करें जो $A$ और $B$ का कोई परिमित अनुक्रम
आगत के रूप में मिलने पर जाँचे कि सभी $A$, सभी $B$ के बाईं ओर हैं या
नहीं। मशीन को आगत अक्षर-शृंखला पट्टी पर छोड़ देनी चाहिए और यदि शृंखला
``वर्णानुक्रम में'' हो तो रुकना चाहिए, अन्यथा सदा चक्रित होना चाहिए।
\end{prob}

\begin{prob}
\emph{वर्णानुक्रमक}: $\{\TMendtape,\TMblank, A, B\}$ वर्णमाला वाली ऐसी
ट्यूरिंग मशीन की अभिकल्पना करें जो $A$ और $B$ के किसी परिमित अनुक्रम को
आगत के रूप में लेकर उन्हें इस प्रकार पुनर्व्यवस्थित करे कि सभी $A$, सभी
~$B$ के बाईं ओर हों। (उदाहरणतः, अनुक्रम $BABAA$ को अनुक्रम $AAABB$ बनना
चाहिए और अनुक्रम $ABBABB$ को अनुक्रम $AABBBB$ बनना चाहिए।)
\end{prob}

\end{document}
